A base UCMF reúne registros de atendimento cardiológico pediátrico. O dicionário descreve a variável \textit{NORMAL X ANORMAL} em termos de presença ou ausência de patologia cardíaca \cite{tema,dict}. Consideramos \textit{Anormal} a classe positiva e investigamos em que medida características antropométricas, demográficas e fisiológicas discriminam esse rótulo. A fonte disponível deixa em aberto a ordem de registro dos campos, a identidade do responsável pela rotulagem e o padrão diagnóstico de referência. Essas lacunas afetam tanto a escolha dos preditores quanto o significado dos resultados.

Na avaliação anterior, o trabalho chegou à transformação dos dados. Nesta A2, revimos seleção, limpeza e transformação e completamos a mineração e a interpretação do processo KDD. O desenho compara famílias de modelos sob as mesmas partições e usa dois conjuntos de atributos para avaliar quanto a conclusão depende de medidas mais próximas da avaliação cardiovascular.
